\section{Agent Tools}
LLMs by themselves are very powerful when it comes to understanding and generating human language.
However, they can be even more powerful when combined with external tools that allow them to perform specific tasks or access up-to-date information. 
In this section, we will explore some of the tools that can be integrated with LLMs to enhance their capabilities.

They enable a wide range of applications, such as:
\begin{itemize}
    \item \textbf{Computational tools:} perform complex calculations, data analysis, and other computational tasks that may be beyond their native capabilities;
    \item \textbf{Knowledge retrieval tools:} access external databases, APIs, or the internet to retrieve up-to-date information;
    \item \textbf{Action execution tools:} interact with other software or hardware to perform actions based on the LLMs output. 
\end{itemize}

\subsection{Register a Tool}
LLM providers usually offers ready-to-use tools, such as a web search tool, code interpreter, file search, and more.
However, they also allow users to create and register custom tools that can be tailored to specific use cases or domains.

To register a tool, we first need to define its schema, which includes the tool's name, description, and input parameters.
An example of a tool schema is shown below:
\begin{lstlisting}[language={}]
    {
    "type": "function",
    "function": {
        "name": "explore_data",
        "description": "Inspect dataset schema, missing values, and candidate targets. Call first when the dataset structure is unknown.",
        "parameters": {
        "properties": {
            "sample_rows": {"type": "integer", "description": "Number of sample rows to return"}
        },
        "required": ["sample_rows"]
        }
    }
    }
\end{lstlisting}

\subsection{Multi-Tool loop}
Real world problems often require the use of multiple tools in a sequence, where the output of one tool serves as the input for another.
The model must reason about which tool to call next based on what it's learned so far.

\paragraph{Example: Data Analysis Agent}
Consider the following example, of a data analysis agent. The following tools are registered:
\begin{itemize}
    \item \textbf{Explore Data:} inspects the dataset schema, missing values, and candidate targets. it has no preconditions;
    \item \textbf{Cluster Data:} performs clustering on the dataset, discovering natural groupings in the data. Needs a numeric column as input;
    \item \textbf{Classify Data:} performs regression or classification on the dataset, depending on the target variable. Needs a target column as input.
\end{itemize}

Different queries may require different sequences of tools. For example:
\begin{itemize}
    \item \textbf{Query 1:} "What kind of data do I have?" $\rightarrow$ Explore Data;
    \item \textbf{Query 2:} "Are there any natural customer segments in the data?" $\rightarrow$ Explore Data $\rightarrow$ Cluster Data;
    \item \textbf{Query 3:} "Segments and churn prediction for age=29, income=50k..." $\rightarrow$ Explore Data $\rightarrow$ Cluster Data $\rightarrow$ Classify Data.
\end{itemize}

\subsection{The Control Problem}
